\scanpage{44}
\mspage{1}
\uline{Groupoïdes de Teichmüller.}

\textcircled{1} Cas topologique.

La variante la plus simple est
\[
\mathcal T^+_{\mathrm{top}}\quad\text{ou}\quad\mathcal T_{\mathrm{top}}
\]
(\uline{groupoïde de Teichmüller \uuline{orienté} topologique})

\uline{Objets} : surfaces \uline{à bord} compactes orientées $(\Sigma)$.

\uline{flèches} : classes d'isotopie d'iso (conservant l'orientation)

\margin{Finalement, \unsure{tout pris en détail}, \unsure{au} point de vue
\unsure{œcuménique}, il semble que c'est le cas $\mathcal T^\pm$ qui offre
le plus \unsure{utilité}; \deleted{\unsure{d'où}} c'est lui qui
\unsure{aura} donc ici la notation la plus simple $\mathcal T$,
et le \unsure{groupe} orienté s'appelle $\mathcal T^+$.}

Variante immédiate
\[
\mathcal T_{\mathrm{top}}^\pm
\]
(\uline{groupoïde de Teichmüller \uuline{biorienté} topologique}) :
mêmes objets, flèches sont les classes d'isotopie d'homéomorphismes
(ne conservant pas néct. l'orientation).

Soit $[\pm1]$ le groupoïde ayant un seul objet $e$ (par exemple $e=\varnothing$)
et comme groupe d'autom $\pm1$. On a un morphisme de signature
\diagram{\begin{tikzcd}[ampersand replacement=\&]
\mathcal T^\pm_{\mathrm{top}}\arrow[r,"\chi"]\&{[\pm1]}
\end{tikzcd}}
dont la \uline{fibre} \deleted{est} (ou noyau) est
$\mathcal T_{\mathrm{top}}^+=\mathcal T_{\mathrm{top}}$.
\margin{\unsure{Ceci n'a un sens} \unsure{qu'}à des surfaces connexes!}

La variante la plus complexe est
\[
S_*\mathcal TD_{\mathrm{top}}^+\quad\text{ou}\quad
S_*\mathcal TD_{\mathrm{top}},
\]
``\uline{groupoïde de Teichmüller \inserted{orienté} spécial dévissé topologique}''
\mspage{43}
